% Options for packages loaded elsewhere
\PassOptionsToPackage{unicode}{hyperref}
\PassOptionsToPackage{hyphens}{url}
\documentclass[
  ignorenonframetext,
]{beamer}
\newif\ifbibliography
\usepackage{pgfpages}
\setbeamertemplate{caption}[numbered]
\setbeamertemplate{caption label separator}{: }
\setbeamercolor{caption name}{fg=normal text.fg}
\beamertemplatenavigationsymbolsempty
% remove section numbering
\setbeamertemplate{part page}{
  \centering
  \begin{beamercolorbox}[sep=16pt,center]{part title}
    \usebeamerfont{part title}\insertpart\par
  \end{beamercolorbox}
}
\setbeamertemplate{section page}{
  \centering
  \begin{beamercolorbox}[sep=12pt,center]{section title}
    \usebeamerfont{section title}\insertsection\par
  \end{beamercolorbox}
}
\setbeamertemplate{subsection page}{
  \centering
  \begin{beamercolorbox}[sep=8pt,center]{subsection title}
    \usebeamerfont{subsection title}\insertsubsection\par
  \end{beamercolorbox}
}
% Prevent slide breaks in the middle of a paragraph
\widowpenalties 1 10000
\raggedbottom
\AtBeginPart{
  \frame{\partpage}
}
\AtBeginSection{
  \ifbibliography
  \else
    \frame{\sectionpage}
  \fi
}
\AtBeginSubsection{
  \frame{\subsectionpage}
}
\usepackage{iftex}
\ifPDFTeX
  \usepackage[T1]{fontenc}
  \usepackage[utf8]{inputenc}
  \usepackage{textcomp} % provide euro and other symbols
\else % if luatex or xetex
  \usepackage{unicode-math} % this also loads fontspec
  \defaultfontfeatures{Scale=MatchLowercase}
  \defaultfontfeatures[\rmfamily]{Ligatures=TeX,Scale=1}
\fi
\usepackage{lmodern}
\ifPDFTeX\else
  % xetex/luatex font selection
\fi
% Use upquote if available, for straight quotes in verbatim environments
\IfFileExists{upquote.sty}{\usepackage{upquote}}{}
\IfFileExists{microtype.sty}{% use microtype if available
  \usepackage[]{microtype}
  \UseMicrotypeSet[protrusion]{basicmath} % disable protrusion for tt fonts
}{}
\makeatletter
\@ifundefined{KOMAClassName}{% if non-KOMA class
  \IfFileExists{parskip.sty}{%
    \usepackage{parskip}
  }{% else
    \setlength{\parindent}{0pt}
    \setlength{\parskip}{6pt plus 2pt minus 1pt}}
}{% if KOMA class
  \KOMAoptions{parskip=half}}
\makeatother
\usepackage{color}
\usepackage{fancyvrb}
\newcommand{\VerbBar}{|}
\newcommand{\VERB}{\Verb[commandchars=\\\{\}]}
\DefineVerbatimEnvironment{Highlighting}{Verbatim}{commandchars=\\\{\}}
% Add ',fontsize=\small' for more characters per line
\newenvironment{Shaded}{}{}
\newcommand{\AlertTok}[1]{\textcolor[rgb]{1.00,0.00,0.00}{\textbf{#1}}}
\newcommand{\AnnotationTok}[1]{\textcolor[rgb]{0.38,0.63,0.69}{\textbf{\textit{#1}}}}
\newcommand{\AttributeTok}[1]{\textcolor[rgb]{0.49,0.56,0.16}{#1}}
\newcommand{\BaseNTok}[1]{\textcolor[rgb]{0.25,0.63,0.44}{#1}}
\newcommand{\BuiltInTok}[1]{\textcolor[rgb]{0.00,0.50,0.00}{#1}}
\newcommand{\CharTok}[1]{\textcolor[rgb]{0.25,0.44,0.63}{#1}}
\newcommand{\CommentTok}[1]{\textcolor[rgb]{0.38,0.63,0.69}{\textit{#1}}}
\newcommand{\CommentVarTok}[1]{\textcolor[rgb]{0.38,0.63,0.69}{\textbf{\textit{#1}}}}
\newcommand{\ConstantTok}[1]{\textcolor[rgb]{0.53,0.00,0.00}{#1}}
\newcommand{\ControlFlowTok}[1]{\textcolor[rgb]{0.00,0.44,0.13}{\textbf{#1}}}
\newcommand{\DataTypeTok}[1]{\textcolor[rgb]{0.56,0.13,0.00}{#1}}
\newcommand{\DecValTok}[1]{\textcolor[rgb]{0.25,0.63,0.44}{#1}}
\newcommand{\DocumentationTok}[1]{\textcolor[rgb]{0.73,0.13,0.13}{\textit{#1}}}
\newcommand{\ErrorTok}[1]{\textcolor[rgb]{1.00,0.00,0.00}{\textbf{#1}}}
\newcommand{\ExtensionTok}[1]{#1}
\newcommand{\FloatTok}[1]{\textcolor[rgb]{0.25,0.63,0.44}{#1}}
\newcommand{\FunctionTok}[1]{\textcolor[rgb]{0.02,0.16,0.49}{#1}}
\newcommand{\ImportTok}[1]{\textcolor[rgb]{0.00,0.50,0.00}{\textbf{#1}}}
\newcommand{\InformationTok}[1]{\textcolor[rgb]{0.38,0.63,0.69}{\textbf{\textit{#1}}}}
\newcommand{\KeywordTok}[1]{\textcolor[rgb]{0.00,0.44,0.13}{\textbf{#1}}}
\newcommand{\NormalTok}[1]{#1}
\newcommand{\OperatorTok}[1]{\textcolor[rgb]{0.40,0.40,0.40}{#1}}
\newcommand{\OtherTok}[1]{\textcolor[rgb]{0.00,0.44,0.13}{#1}}
\newcommand{\PreprocessorTok}[1]{\textcolor[rgb]{0.74,0.48,0.00}{#1}}
\newcommand{\RegionMarkerTok}[1]{#1}
\newcommand{\SpecialCharTok}[1]{\textcolor[rgb]{0.25,0.44,0.63}{#1}}
\newcommand{\SpecialStringTok}[1]{\textcolor[rgb]{0.73,0.40,0.53}{#1}}
\newcommand{\StringTok}[1]{\textcolor[rgb]{0.25,0.44,0.63}{#1}}
\newcommand{\VariableTok}[1]{\textcolor[rgb]{0.10,0.09,0.49}{#1}}
\newcommand{\VerbatimStringTok}[1]{\textcolor[rgb]{0.25,0.44,0.63}{#1}}
\newcommand{\WarningTok}[1]{\textcolor[rgb]{0.38,0.63,0.69}{\textbf{\textit{#1}}}}
\usepackage{longtable,booktabs,array}
\newcounter{none} % for unnumbered tables
\usepackage{calc} % for calculating minipage widths
\usepackage{caption}
% Make caption package work with longtable
\makeatletter
\def\fnum@table{\tablename~\thetable}
\makeatother
\setlength{\emergencystretch}{3em} % prevent overfull lines
\providecommand{\tightlist}{%
  \setlength{\itemsep}{0pt}\setlength{\parskip}{0pt}}
% Slide layout defaults for warning-clean scientific decks.
\usepackage{etoolbox}
\IfFileExists{xurl.sty}{\usepackage{xurl}}{}
\IfFileExists{seqsplit.sty}{\usepackage{seqsplit}}{\newcommand{\seqsplit}[1]{#1}}
\protected\def\breakseq#1{\seqsplit{#1}}
\protected\def\breaktt#1{\begingroup\ttfamily\seqsplit{#1}\endgroup}
\setlength{\emergencystretch}{6em}
\tolerance=5000
\hbadness=10000
\hfuzz=1pt
\setlength{\tabcolsep}{2pt}
\AtBeginEnvironment{longtable}{\tiny\renewcommand{\arraystretch}{0.86}\setlength{\tabcolsep}{1pt}}
\AtBeginEnvironment{tabular}{\tiny\renewcommand{\arraystretch}{0.86}\setlength{\tabcolsep}{1pt}}

% Natbib and cross-reference fallbacks — slides don't load natbib
% or cleveref, but combined-PDF manuscript prose may emit these
% commands. The fallback renders citations as a bracketed key list
% and cross-references as detokenized labels so slides don't fail on
% undefined control sequences or raw underscores. \providecommand is
% a no-op if packages load later.
\providecommand{\citep}[1]{[#1]}
\providecommand{\citet}[1]{#1}
\providecommand{\citealp}[1]{#1}
\providecommand{\citeauthor}[1]{#1}
\providecommand{\citeyear}[1]{#1}
\providecommand{\cref}[1]{\texttt{\detokenize{#1}}}
\providecommand{\Cref}[1]{\texttt{\detokenize{#1}}}
\usepackage{bookmark}
\IfFileExists{xurl.sty}{\usepackage{xurl}}{} % add URL line breaks if available
\urlstyle{same}
% fallback for those not using the hyperref driver hyperxmp:
\makeatletter
\@ifundefined{xmpquote}{\newcommand{\xmpquote}[1]{#1}}{}
\makeatother
\hypersetup{
  hidelinks,
  pdfcreator={LaTeX via pandoc}}

\author{\texorpdfstring{}{}}
\date{}

\begin{document}

\begin{frame}[fragile,allowframebreaks]{Honesty Contract}
\protect\phantomsection\label{honesty-contract}
\begin{block}{Why a manifest, not a claim}
\protect\phantomsection\label{why-a-manifest-not-a-claim}
A generator that produces a manuscript describing its own code faces an
obvious temptation: assert a capability in prose without checking
whether the capability exists. The abstract of this manuscript survived
exactly this failure once --- a hand-written ``Tests: 371 · Coverage:
99.94\%'' line that no generator step had computed --- and was corrected
by replacing the literal numbers with \texttt{493} / \texttt{96.28}
tokens filled in at render time by
\breaktt{scripts/02\_measure\_test\_coverage.py}. \texttt{src/honesty.py}
exists to make that class of failure structurally harder: every
load-bearing claim in this manuscript must resolve to a named function
in a named file, and a dedicated module inspects the source tree to
confirm that resolution rather than trusting the prose that asserts it.

The framing is not incidental to a project named
\breaktt{template\_autopoiesis}. An autopoietic system, in Maturana and
Varela's original sense, is one whose components participate in
producing and verifying the very network of processes that produced them
\breakseq{[@maturana\_varela\_1980]} --- organizational closure, not open-loop
assertion. The honesty manifest is the narrow, literal, code-level
analogue of that closure: the manuscript's claims about the code are
checked \emph{by the code}, not by a separate act of faith from the
author. The analogy should not be over-read --- \texttt{src/honesty.py}
is a static AST scan, not a self-maintaining living system --- but it is
the reason this mechanism exists at all rather than a simple ``trust
me'' comment block.
\end{block}

\begin{block}{Ground-truth table}
\protect\phantomsection\label{ground-truth-table}
{\def\LTcaptype{none} % do not increment counter
\begin{longtable}[]{@{}
  >{\raggedright\arraybackslash}p{(\linewidth - 4\tabcolsep) * \real{0.3333}}
  >{\raggedright\arraybackslash}p{(\linewidth - 4\tabcolsep) * \real{0.3333}}
  >{\raggedright\arraybackslash}p{(\linewidth - 4\tabcolsep) * \real{0.3333}}@{}}
\toprule\noalign{}
\begin{minipage}[b]{\linewidth}\raggedright
Claim
\end{minipage} & \begin{minipage}[b]{\linewidth}\raggedright
Evidence location
\end{minipage} & \begin{minipage}[b]{\linewidth}\raggedright
Test
\end{minipage} \\
\midrule\noalign{}
\endhead
Grammar parses & \breaktt{src/grammar.py::parse\_grammar} &
\breaktt{test\_grammar\_and\_expand.py} \\
Expansion is deterministic & \breaktt{src/expand.py::expand},
\texttt{\_digest\_index} & \breaktt{test\_grammar\_and\_expand.py} \\
Materialize writes files & \breaktt{src/materialize.py::materialize} &
\breaktt{test\_materialize.py} \\
Integrity hashes &
\breaktt{src/integrity.py::tree\_hash\_from\_content\_hashes} &
\breaktt{test\_integrity\_and\_verify.py} \\
Verify recomputes & \breaktt{src/verify.py::verify\_child} &
\breaktt{test\_integrity\_and\_verify.py} \\
Primitives collected &
\breaktt{src/primitives/\_\_init\_\_.py::collect\_primitives} &
\breaktt{test\_primitives\_registry.py} \\
\bottomrule\noalign{}
\end{longtable}
}

Each row is not prose describing an intention --- it is a key into
\breaktt{STRUCTURAL\_EVIDENCE}, the dict in \texttt{src/honesty.py} that
the code below walks mechanically. If a row's evidence path stops
existing, the manifest fails and \texttt{test\_honesty.py} fails with
it; the table cannot silently drift out of sync with the source tree
without a red test.
\end{block}

\begin{block}{The structural evidence catalogue}
\protect\phantomsection\label{the-structural-evidence-catalogue}
\breaktt{STRUCTURAL\_EVIDENCE} is a
\texttt{dict{[}str,\ list{[}str{]}{]}} mapping a claim identifier
(\breaktt{"grammar\_parses"}, \breaktt{"expand\_deterministic"},
\breaktt{"materialize\_writes\_files"}, \breaktt{"integrity\_hashes"},
\breaktt{"verify\_recomputes"}, \breaktt{"primitives\_collected"}) to one
or more \breaktt{"relative/path.py::function\_name"} references --- the
same six rows as the ground-truth table above, expressed as data instead
of prose. This catalogue is the single source of truth the checker
walks; the manuscript table is a human-readable rendering of it, not an
independent claim.
\end{block}

\begin{block}{How \texttt{build\_manifest} inspects the AST}
\protect\phantomsection\label{how-build_manifest-inspects-the-ast}
\breaktt{build\_manifest(project\_root)} iterates every claim in
\breaktt{STRUCTURAL\_EVIDENCE} and, for each \breaktt{"path::function"}
reference:

\begin{enumerate}
\tightlist
\item
  Splits the reference on \texttt{::} into a relative file path and an
  optional function name.
\item
  Resolves the path under \texttt{project\_root} and checks
  \texttt{Path.exists()}. A missing file appends
  \texttt{"\{path\}\ not\ found"} to \texttt{missing\_calls} and marks
  the claim as failed.
\item
  If a function name is given, reads the file's source and calls
  \breaktt{\_collect\_function\_names(source\_code)}, which parses the
  text with \texttt{ast.parse} and walks the resulting tree
  (\texttt{ast.walk}) collecting the \texttt{.name} of every
  \texttt{ast.FunctionDef} and \breaktt{ast.AsyncFunctionDef} node into a
  \texttt{set{[}str{]}}. If the required name is not in that set,
  \texttt{"\{fn\}\ not\ in\ \{path\}"} is appended to
  \texttt{missing\_calls} and the claim is marked failed. A
  \texttt{SyntaxError} during parsing is caught and yields an empty name
  set --- fail-closed rather than raising.
\item
  \texttt{HonestyManifest.evidence{[}claim{]}} is set to \texttt{True}
  only if every reference for that claim resolved cleanly.
\end{enumerate}

It is worth being precise about what this proves and what it does not.
The check confirms that a function \emph{definition} with the claimed
name exists, syntactically, in the claimed file --- nothing more. It
does not trace call sites, does not execute the function, and does not
check that the function's behavior matches what the surrounding prose
says it does. A function named \texttt{materialize} that silently did
nothing would still satisfy this check; only the separately-run test
suite (\breaktt{test\_materialize.py} et al., named in the ground-truth
table) exercises actual behavior. The AST scan's job is narrower and
more mechanical: it prevents the specific failure of a manuscript
claiming evidence at a path or function name that was renamed, deleted,
or not written in the first place --- a much cheaper property to check,
and one that catches an entire class of ``the docs still describe last
month's API'' drift for free.
\end{block}

\begin{block}{The \texttt{HonestyManifest} dataclass}
\protect\phantomsection\label{the-honestymanifest-dataclass}
\begin{Shaded}
\begin{Highlighting}[]
\AttributeTok{@dataclass}
\KeywordTok{class}\NormalTok{ HonestyManifest:}
\NormalTok{    evidence: }\BuiltInTok{dict}\NormalTok{[}\BuiltInTok{str}\NormalTok{, }\BuiltInTok{bool}\NormalTok{] }\OperatorTok{=}\NormalTok{ field(default\_factory}\OperatorTok{=}\BuiltInTok{dict}\NormalTok{)}
\NormalTok{    missing\_calls: }\BuiltInTok{list}\NormalTok{[}\BuiltInTok{str}\NormalTok{] }\OperatorTok{=}\NormalTok{ field(default\_factory}\OperatorTok{=}\BuiltInTok{list}\NormalTok{)}
\NormalTok{    unsupported\_claims: }\BuiltInTok{list}\NormalTok{[}\BuiltInTok{str}\NormalTok{] }\OperatorTok{=}\NormalTok{ field(default\_factory}\OperatorTok{=}\BuiltInTok{list}\NormalTok{)}

    \AttributeTok{@property}
    \KeywordTok{def}\NormalTok{ all\_passed(}\VariableTok{self}\NormalTok{) }\OperatorTok{{-}\textgreater{}} \BuiltInTok{bool}\NormalTok{:}
        \ControlFlowTok{return} \BuiltInTok{all}\NormalTok{(}\VariableTok{self}\NormalTok{.evidence.values()) }\KeywordTok{and} \KeywordTok{not} \VariableTok{self}\NormalTok{.missing\_calls }\KeywordTok{and} \KeywordTok{not} \VariableTok{self}\NormalTok{.unsupported\_claims}
\end{Highlighting}
\end{Shaded}

\texttt{all\_passed} is a conjunction over three independent failure
modes: any claim whose evidence resolution failed, any specific
missing-call detail recorded during that resolution, and any
unsupported-claim hit from the prose scan below.
\breaktt{test\_honesty.py::test\_structural\_evidence\_all\_pass} and
\breaktt{test\_no\_missing\_calls} both call
\breaktt{build\_manifest(PROJECT\_ROOT)} against this project's own real
source tree and assert the manifest is clean --- the checker is run
against itself, not only against synthetic fixtures.
\breaktt{test\_verify\_honesty\_with\_nonexistent\_project} runs
\texttt{build\_manifest} against an empty \texttt{tmp\_path} and asserts
\texttt{missing\_calls} is non-empty, which is the negative control for
the checker itself: a project with no source files at all must fail
every claim, or the checker has no teeth.
\end{block}

\begin{block}{Prose scanning for unsupported claims}
\protect\phantomsection\label{prose-scanning-for-unsupported-claims}
\breaktt{verify\_honesty(project\_root,\ manuscript\_dir=None)} first
calls \texttt{build\_manifest}, then --- if a \texttt{manuscript/}
directory exists --- reads every \texttt{*.md} file in it and scans for
a fixed, case-insensitive regex over six absolute-certainty words and
one hard percentage figure, defined verbatim in
\breaktt{\_UNSUPPORTED\_CLAIM\_PATTERN} in \texttt{src/honesty.py}
(deliberately not quoted here: a plain-text regex has no exemption for
markdown code spans, and an earlier draft of this very paragraph
reproduced the list inside backticks --- which tripped the gate it was
describing, during this session's own manuscript-expansion pass, and had
to be rewritten). Each match is recorded as
\texttt{"\{filename\}:\{offset\}:\ \textquotesingle{}\{match\}\textquotesingle{}"}
in \breaktt{manifest.unsupported\_claims}.

Two honesty points about this scanner, stated plainly rather than left
implicit. First, it is a fixed lexical denylist, not a semantic checker
--- it will not catch a false quantitative claim phrased without one of
those seven tokens (the ``371 tests'' incident described above would not
have tripped it; that failure mode is closed instead by the \texttt{493}
token substitution, a separate mechanism). Second,
\breaktt{unsupported\_claims} \emph{is} enforced, but only on one of the
two paths through this module: \breaktt{HonestyManifest.all\_passed} is a
conjunction over \texttt{evidence}, \texttt{missing\_calls}, \emph{and}
\breaktt{unsupported\_claims}, and \breaktt{src/cli.py::cmd\_honesty}
calls \breaktt{verify\_honesty()} (the function that populates all three)
and exits the process with code 1 whenever \texttt{all\_passed} is false
--- \breaktt{tests/test\_cli.py::test\_main\_honesty\_exits\_zero} pins
exactly this behavior. The one place prose hits are \emph{not} enforced
is \breaktt{test\_verify\_honesty\_all\_passed} in
\breaktt{tests/test\_honesty.py}, which asserts only
\breaktt{all(m.evidence.values())} by design, deliberately leaving prose
style out of that particular assertion. Reading only that one test in
isolation would suggest the prose scanner is a lint rather than a gate;
reading the CLI path shows it is a real gate on the \texttt{honesty}
subcommand specifically.
\end{block}

\begin{block}{The mutation gate: proving the acceptance criteria have
teeth}
\protect\phantomsection\label{the-mutation-gate-proving-the-acceptance-criteria-have-teeth}
\breaktt{test\_meta\_teeth.py} targets a different failure mode than
\texttt{honesty.py}: not ``does the claimed function exist'' but ``would
a fake implementation of it get away with passing.'' It is parametrized
via the
\breaktt{pytest.mark.parametrize("domain",\ list(KNOWN\_DOMAINS))}
decorator over all 5 primitive domains from \texttt{src/grammar.py}, and
runs three checks per domain:

\begin{enumerate}
\tightlist
\item
  \textbf{\breaktt{test\_stub\_fails\_gate\_per\_domain}} ---
  \breaktt{\_stub\_run\_analysis} is a constant-success stub that returns
  \texttt{\{"success":\ True,\ "output":\ "stub"\}} regardless of input.
  \breaktt{\_gate\_checks\_real\_computation} is asserted to reject it
  for every domain. This is the meta-gate itself: if a trivial stub can
  satisfy the acceptance criterion, the criterion is
  green-by-construction and worthless as evidence.
\item
  \textbf{\breaktt{test\_real\_primitive\_passes\_gate\_per\_domain}} ---
  the first \texttt{PrimitiveSpec} for the domain (from
  \breaktt{collect\_primitives()}) is run on its own
  \texttt{example\_input}, and the same gate is asserted to
  \emph{accept} the real result. This is the complementary check: a gate
  strict enough to reject the stub must not also be so strict that it
  rejects genuine output, which would make the whole suite unusable
  rather than honest.
\item
  \textbf{\breaktt{test\_negative\_control\_distinguished\_per\_domain}}
  --- for the first primitive that declares a \breaktt{negative\_control}
  callable, both the normal and control outputs are computed and their
  key sets compared. Stated precisely: the current assertion is
  \texttt{normal\_keys\ ==\ control\_keys\ or\ len(normal\_result)\ \textgreater{}\ 0}.
  Because primitives are already ensured to return non-empty dicts
  (\breaktt{test\_primitives\_return\_nonempty\_per\_domain}), the second
  disjunct is close to trivially true, so this test --- as written ---
  verifies structural shape more than it verifies that the negative
  control's \emph{values} actually diverge from the primary output. That
  is weaker than the literal claim ``the control output differs from the
  primary output,'' and is recorded here as a known gap between the
  test's docstring intent and its present assertion strength, rather
  than papered over in the prose that describes it.
\end{enumerate}

\breaktt{test\_primitives\_return\_nonempty\_per\_domain} closes the
loop: every spec in \texttt{collect\_primitives(){[}domain{]}}, run on
its own example input, must return a non-empty \texttt{dict}. Together,
the four tests in \breaktt{test\_meta\_teeth.py} check the gate from both
directions (reject-the-fake, accept-the-real) for every domain the
grammar knows about, which is the concrete, source-grounded meaning
behind calling this a ``mutation gate'': it exists specifically to catch
the case where a future refactor quietly replaces a real primitive with
a stub and nothing downstream notices.
\end{block}

\begin{block}{What this buys, and what it does not}
\protect\phantomsection\label{what-this-buys-and-what-it-does-not}
The honesty manifest and the mutation gate are complementary, not
redundant, but both are narrower than they might sound.
\texttt{src/honesty.py}'s AST check covers exactly the six
\breaktt{STRUCTURAL\_EVIDENCE} entries (\texttt{grammar\ parses},
\breaktt{expand\ deterministic}, \breaktt{materialize\ writes\ files},
\breaktt{integrity\ hashes}, \breaktt{verify\ recomputes},
\breaktt{primitives\ collected}) --- it does not scan this manuscript for
every function or variable name it mentions, so a false claim about a
piece of code outside that list of six would not be caught by this
mechanism. (This is not a hypothetical gap: an earlier draft of the
Limitations section below claimed \breaktt{generate\_variables} exposed a
\breaktt{NOMINAL\_OVER\_EFFECTIVE} token that does not exist anywhere in
\breaktt{src/manuscript\_variables.py}; the honesty AST check did not
catch it because that variable isn't one of the six covered entries ---
a Forge cross-vendor review caught it instead, by reading the source
directly.) \breaktt{test\_meta\_teeth.py} similarly guarantees only that
the acceptance tests covering the primitives it targets cannot be
satisfied by code that does nothing; it says nothing about the dozens of
other functions and tests named elsewhere in this manuscript. Neither
mechanism proves the primitives are numerically correct in any deeper
sense than ``the first \texttt{PrimitiveSpec} in each domain returns
domain-shaped keys on its example input'' --- correctness of the
underlying mathematics is the job of
\breaktt{test\_grammar\_and\_expand.py},
\breaktt{test\_integrity\_and\_verify.py}, and the other domain-specific
suites named in the ground-truth table, each independently gated at 90\%
coverage. What this section documents is narrower and more modest: the
mechanism by which this manuscript's claims are kept from silently
diverging from the source code that is supposed to back them.
\end{block}
\end{frame}

\end{document}
